%%% FIELD clipping-transport-dual-statement
For every finite signed \(\sigma\) on \(S\), \[H_{M_d}(\sigma)=\min_{\pi\in\Pi_{\le}(\sigma)}\left\{\sigma^+(S)-\pi(S\times S)+M_d\int_{S\times S}\|x-y\|_2\,d\pi(x,y)\right\},\] and the minimum is attained. After the affine transformation \(\phi=2g_d-1\) and the corresponding ground-metric rescaling, this support identity is equivalent to Piccoli and Rossi's generalized \(W_1^{1,1}\) flat-metric dual. For every \(m_d\in K_d\), \[U_d(m_d)=\inf_{\lambda_d\in\mathbb R^{|J_d|}}D_d(\lambda_d;m_d),\qquad L_d(m_d)=1-U_d(\mathbf 1-m_d).\] No condition \(\sum_{j\in J_d}\lambda_{dj}=1\) is imposed. A finite minimizing \(\lambda_d\) exists whenever \(r_d(m_d)>0\). In the one-cell unit-rectangle specialization with uniform \(\mu_{dj}\), uniform boundary \(\nu\), \(M_d=1\), and \(m_d=0\), the primal value is zero but every finite multiplier has strictly positive dual objective, so multiplier attainment fails at that boundary mean.
%%% FIELD clipping-transport-dual-proof
Fix an arm \(d\), and abbreviate \(M=M_d\), \(\mathcal F=\mathcal F_d\), \(A=A_d\), \(m=m_d\), and
\[
 T(g)=\int_B g\,d\nu.
\]
The range and Lipschitz assumptions say, separately, that every \(g\in\mathcal F\) takes values in \([0,1]\) and satisfies \(\lvert g(x)-g(y)\rvert\leq M\lVert x-y\rVert_2\).  The polygonal-support assumption makes \(S\) compact and convex, hence connected.  The measures \(\mu_{dj}\) and \(\nu\) are probability measures supported on \(S\), so every integral below is finite.

We first record the compactness needed for all maxima.  Given a sequence \((g_n)\) in \(\mathcal F\), choose, for each positive integer \(k\), a finite \(k^{-1}\)-net of \(S\).  Successive subsequence extraction and a diagonal argument produce a subsequence whose values converge at every point in the union of these nets.  If \(M>0\), then, for any \(\varepsilon>0\), a sufficiently fine finite net and the common Lipschitz bound give, for all sufficiently large indices \(n,n'\),
\[
 \sup_{x\in S}\lvert g_n(x)-g_{n'}(x)\rvert
 \leq 2M\delta+\max_{z\text{ in the finite net}}
       \lvert g_n(z)-g_{n'}(z)\rvert
 <\varepsilon.
\]
If \(M=0\), connectedness makes every \(g_n\) constant and the same conclusion follows from compactness of \([0,1]\).  Thus a subsequence converges uniformly.  Its limit remains \([0,1]\)-valued and \(M\)-Lipschitz, so \(\mathcal F\) is compact in the uniform norm; it is convex by the two defining inequalities.  For every finite signed measure \(\zeta\),
\[
 \left\lvert\int_S(g-\widetilde g)\,d\zeta\right\rvert
 \leq \lVert g-\widetilde g\rVert_\infty\lvert\zeta\rvert(S),
\]
so integration is continuous.  In particular, the maximum in the support-cost definition is attained.  The same compactness also implies attainment of the arm endpoint maxima and minima whenever their exact-moment fiber is nonempty; mean compatibility supplies that nonemptiness.

We now derive the clipping-aware partial-transport formula.  Let \(\sigma=\sigma^+-\sigma^-\) be the Jordan decomposition and put
\[
 p=\sigma^+(S),\qquad q=\sigma^-(S).
\]
Suppose first that \(M>0\), define the dilation \(R(x)=2Mx\), and let
\[
 \alpha=R_\#\sigma^+,
 \qquad
 \beta=R_\#\sigma^-.
\]
For \(g\in\mathcal F\), define \(\phi(Rx)=2g(x)-1\).  Then
\[
 \lVert\phi\rVert_\infty\leq1,
 \qquad
 \lvert\phi(Rx)-\phi(Ry)\rvert
 \leq 2M\lVert x-y\rVert_2
 =\lVert Rx-Ry\rVert_2.
\]
Conversely, any function with these two bounds on \(R(S)\) gives a member of \(\mathcal F\) through \(g(x)=\{1+\phi(Rx)\}/2\).  There is no extension gap between functions on \(R(S)\) and functions on the ambient Euclidean space: the explicit formula
\[
 \Phi(z)=\inf_{u\in R(S)}\{\phi(u)+\lVert z-u\rVert_2\}
\]
is one-Lipschitz and equals \(\phi\) on \(R(S)\), while clipping \(\Phi\) to \([-1,1]\) preserves both that equality and the one-Lipschitz bound.  Therefore the bounded-Lipschitz supremum on \(R(S)\) is the ambient flat-metric supremum.

The cited Piccoli--Rossi flat-duality result, \cite[Definition 12 and Theorem 13]{PiccoliRossi2016GeneralizedWasserstein}, now gives
\[
\begin{aligned}
 H_M(\sigma)
 &=\max_{g\in\mathcal F}\int_S g\,d\sigma\\
 &=\frac{p-q}{2}
   +\frac12\sup_{\substack{\lVert\phi\rVert_\infty\leq1\\
                             \operatorname{Lip}(\phi)\leq1}}
       \int \phi\,d(\alpha-\beta)\\
 &=\frac{p-q}{2}+\frac12W_1^{1,1}(\alpha,\beta).
\end{aligned}
\]
This proves the stated affine and ground-metric equivalence, rather than merely an analogy with the generalized distance.

Apply next the cited attained submeasure representation, \cite[Definition 3 and Proposition 4]{PiccoliRossi2016GeneralizedWasserstein}.  If \(\widetilde\alpha\leq\alpha\) and \(\widetilde\beta\leq\beta\) have common mass \(t\), and \(\overline\pi\) couples them, pull \(\overline\pi\) back through \(R\times R\) to obtain \(\pi\).  Then
\[
 \pi_1\leq\sigma^+,
 \qquad
 \pi_2\leq\sigma^-,
 \qquad
 \pi(S\times S)=t,
\]
and
\[
 \int\lVert u-v\rVert_2\,d\overline\pi(u,v)
 =2M\int_{S\times S}\lVert x-y\rVert_2\,d\pi(x,y).
\]
Conversely, every \(\pi\in\Pi_{\leq}(\sigma)\) pushes forward to such an equal-mass pair of submeasures and a coupling of them.  Substituting the deletion costs \(p-t\) and \(q-t\) into the two cited identities therefore gives, equation by equation,
\[
\begin{aligned}
 H_M(\sigma)
 &=\min_{\pi\in\Pi_{\leq}(\sigma)}
   \left\{
      \frac{p-q}{2}+\frac{p+q-2\pi(S\times S)}{2}
      +M\int_{S\times S}\lVert x-y\rVert_2\,d\pi(x,y)
   \right\}\\
 &=\min_{\pi\in\Pi_{\leq}(\sigma)}
   \left\{
      p-\pi(S\times S)
      +M\int_{S\times S}\lVert x-y\rVert_2\,d\pi(x,y)
   \right\}.
\end{aligned}
\]
The minimizing submeasures in the cited representation are attained.  For those fixed submeasures, a minimizing transport coupling is also attained: a minimizing sequence has fixed finite mass on the compact set \(R(S)\times R(S)\); diagonal convergence of its integrals against a countable uniformly dense family of continuous functions gives a positive limit measure, the two marginals pass to the limit, and the continuous distance cost converges.  Pullback therefore gives an attaining \(\pi\in\Pi_{\leq}(\sigma)\).

If \(M=0\), every class member is a constant \(c\in[0,1]\), because \(S\) is connected.  Hence
\[
 H_0(\sigma)=\max_{0\leq c\leq1}c(p-q)=(p-q)_+.
\]
Every partial coupling has mass at most \(\min\{p,q\}\).  That mass is attained by the zero plan when \(pq=0\), and otherwise by
\[
 \pi=\min\{p,q\}\left(\frac{\sigma^+}{p}\right)
                    \otimes\left(\frac{\sigma^-}{q}\right).
\]
Thus the transport expression with zero ground cost equals
\[
 p-\min\{p,q\}=(p-q)_+,
\]
and its minimum is attained.  This completes the support identity for every \(M\geq0\), without any equality-of-total-mass condition on \(\sigma\).

We next dualize the exact areal moments.  Define the compact convex set
\[
 \mathcal C=\{(Ag,T(g)):g\in\mathcal F\}
 \subseteq \mathbb R^{\lvert J_d\rvert}\times\mathbb R.
\]
Compactness follows from the compactness of \(\mathcal F\) and continuity of the finitely many probability-measure integrals; convexity follows from linearity.  Let \(u=U_d(m)\).  By mean compatibility and the endpoint definition, there is \(g_0\in\mathcal F\) such that
\[
 Ag_0=m,
 \qquad
 T(g_0)=u.
\]
For any unrestricted \(\lambda\in\mathbb R^{\lvert J_d\rvert}\), the support-cost and dual-objective definitions yield
\[
\begin{aligned}
 D_d(\lambda;m)
 &=\lambda^\top m+\sup_{g\in\mathcal F}
       \left\{T(g)-\lambda^\top Ag\right\}\\
 &=\sup_{(v,t)\in\mathcal C}
       \left\{t-\lambda^\top(v-m)\right\}.
\end{aligned}
\]
Every feasible \(g\) with \(Ag=m\) consequently satisfies \(T(g)\leq D_d(\lambda;m)\).  Maximizing over the fiber and then infimizing over \(\lambda\) proves weak duality,
\[
 u\leq\inf_{\lambda\in\mathbb R^{\lvert J_d\rvert}}D_d(\lambda;m).
\]

For the reverse inequality, let \(E=\operatorname{span}(K_d-K_d)\) and work in the finite-dimensional affine Euclidean space
\[
 \mathcal L=(m+E)\times\mathbb R,
\]
which contains \(\mathcal C\).  For \(\eta>0\), the point \(y_\eta=(m,u+\eta)\) is not in \(\mathcal C\), because \(u\) is the greatest second coordinate in the fiber over \(m\).  Choose a nearest point \(z_\eta\in\mathcal C\) and write the unit vector from \(z_\eta\) to \(y_\eta\) as
\[
 n_\eta=(a_\eta,b_\eta)\in E\times\mathbb R.
\]
For every \(c\in\mathcal C\), comparison of \(z_\eta\) with \(z_\eta+s(c-z_\eta)\), followed by differentiation of squared distance at \(s=0\), gives
\[
 n_\eta^\top(c-z_\eta)\leq0.
\]
Since \(n_\eta^\top(y_\eta-z_\eta)>0\), comparison with \((m,u)\in\mathcal C\) gives
\[
 b_\eta\eta=n_\eta^\top\{y_\eta-(m,u)\}>0,
\]
so \(b_\eta>0\).  For \(c=(v,t)\in\mathcal C\), the same separating inequality implies
\[
 a_\eta^\top(v-m)+b_\eta(t-u-\eta)\leq0.
\]
Set \(\lambda_\eta=-a_\eta/b_\eta\), viewed as an unrestricted vector in \(\mathbb R^{\lvert J_d\rvert}\).  Then
\[
 t-\lambda_\eta^\top(v-m)\leq u+\eta
 \quad\text{for every }(v,t)\in\mathcal C,
\]
and hence
\[
 D_d(\lambda_\eta;m)\leq u+\eta.
\]
Letting \(\eta\downarrow0\) and combining this inequality with weak duality proves
\[
 U_d(m)=\inf_{\lambda\in\mathbb R^{\lvert J_d\rvert}}D_d(\lambda;m).
\]
The separating vector ranged over the whole multiplier space; neither the argument nor the dual-objective definition imposes \(\sum_{j\in J_d}\lambda_{dj}=1\).

Suppose now that \(r_d(m)>0\).  By the relative-interior-radius definition, there is a \(\rho>0\) such that
\[
 m+\rho\mathbb B_E\subseteq K_d.
\]
Apply the preceding normalized separation with \(\eta=1/n\).  A subsequence of the unit normals \((a_n,b_n)\) converges to some unit vector \((a,b)\in E\times\mathbb R\).  Passing to the limit in the separating inequality gives
\[
 a^\top(v-m)+b(t-u)\leq0
 \quad\text{for every }(v,t)\in\mathcal C.
\]
The limit cannot have \(b=0\).  Indeed, then \(\lVert a\rVert_2=1\), while \(m+\rho a\in K_d\); choosing \(g\in\mathcal F\) with \(Ag=m+\rho a\) would give \(\rho=a^\top(Ag-m)\leq0\), a contradiction.  Therefore \(b>0\).  With \(\lambda=-a/b\),
\[
 t-\lambda^\top(v-m)\leq u
 \quad\text{for every }(v,t)\in\mathcal C,
\]
so \(D_d(\lambda;m)\leq u\).  Weak duality gives the reverse inequality, and hence this finite \(\lambda\) attains the dual minimum.

For the lower endpoint, probability normalization of every \(\mu_{dj}\) and of \(\nu\) makes \(g\mapsto1-g\) a bijection of \(\mathcal F\) satisfying
\[
 A(1-g)=\mathbf 1-Ag,
 \qquad
 T(1-g)=1-T(g).
\]
In particular, \(m\in K_d\) implies \(\mathbf1-m\in K_d\), and reflection of the exact-moment fiber yields
\[
 L_d(m)=1-U_d(\mathbf1-m).
\]

Finally, specialize to the one-cell rectangle
\[
 S=C_{dj}=[0,1]\times[0,1],
 \qquad
 B=\{0\}\times[0,1],
\]
with uniform area law \(\mu_{dj}\), uniform length law \(\nu\), \(M=1\), and \(m=0\).  If \(g\) is feasible, then \(g\geq0\) and \(\int_Sg\,d\mu_{dj}=0\).  Were \(g\) positive at any point, continuity would make it positive on a relative neighborhood of positive uniform area, contradicting the zero integral.  Thus \(g\equiv0\), including on \(B\), and
\[
 U_d(0)=0.
\]
For any finite scalar multiplier \(\lambda\) and any \(t\in(0,1]\), the cone
\[
 g_t(x,y)=(t-x)_+
\]
is \([0,1]\)-valued and one-Lipschitz, with
\[
 \int_Bg_t\,d\nu=t,
 \qquad
 \int_Sg_t\,d\mu_{dj}
 =\int_0^t(t-x)\,dx
 =\frac{t^2}{2}.
\]
The support definition therefore implies
\[
 D_d(\lambda;0)=H_1(\nu-\lambda\mu_{dj})
 \geq t-\frac{\lambda t^2}{2}.
\]
If \(\lambda\leq0\), take any \(t\in(0,1]\); if \(\lambda>0\), take \(0<t<\min\{1,2/\lambda\}\).  In either case the right-hand side is strictly positive.  Thus every finite multiplier has \(D_d(\lambda;0)>0\), although the proved value identity gives \(\inf_\lambda D_d(\lambda;0)=U_d(0)=0\).  No finite multiplier attains the infimum at this boundary mean, whereas the primal surface and the partial-transport minimizer do attain their respective optima.
